\subsection{Three-row ordinary acts through thirteen corner elements}

\begin{proposition}[Proved: the four forced faithful targets]
\label{CP2:A-forced-three-row-hulls-thirteen}
Among labelled faithful actor images with minimal band $3\times2$,
balanced row values $(0,0,1)$, order at most thirteen, and all four
profiles
\[
 p_i=((i,i,2),(0,1)),\qquad
 r_i=((2,2,i),(1,0)),\qquad i=0,1,
\]
there are exactly four targets. They have orders $11,12,13,13$.
Their full row augmentation modules have projective dimension
at most one over their contracted rational algebras.
\end{proposition}
\begin{proof}
The balanced actor monoid has the24 elements described in
Lemma~\ref{CP2:A-three-by-two-balanced-actor-cover}. Its six band
actors, identity and the four displayed actors form an11-element
submonoid. There are thirteen remaining actors, and any target
in the stated scope adds at most two. Trying all
$1+13+\binom{13}{2}=92$ subsets and checking every product
therefore has exact labelled scope. The only closed choices are
the empty choice, the singleton
\[
 s=((1,0,2),(0,1)),
\]
and the two pairs, indexed by $\epsilon=0,1$,
\[
 x=((0,0,1),(\epsilon,\epsilon)),\qquad
 y=((1,1,0),(1-\epsilon,1-\epsilon)).
\]
The independent checker also obtains these same four sets by
successively adjoining an actor and taking full multiplication
closure, pruning only a set larger than thirteen. Any target
can be generated in this way without ever leaving itself.
No ordinary-isomorphism reduction is used in either enumeration.

Here are exact projective presentations for all four targets.
In the contracted algebra $A$, put
\[
 e_0=(p_0-r_0)/2,\qquad e_1=(p_1-r_1)/2,
 \qquad f=1-p_0.
\]
Both $e_0$ and $f$ are idempotent. Let the row module $V$ have
basis $u=[1]-[0]$, $v=[2]-[0]$.
For the11- and12-element targets,
\[
 Ae_0=\mathbb Qe_0\oplus\mathbb Qe_1,\qquad
 e_0v=v,\qquad e_1v=v-u.
\]
Thus multiplication on $v$ identifies the projective ideal
$Ae_0$ with $V$.

For either13-element target put $z=(x-y)/2$. Direct multiplication
gives
\[
 Ae_0=\mathbb Qe_0\oplus\mathbb Qe_1\oplus\mathbb Qz,
 \qquad zv=u.
\]
The same map $\phi:Ae_0\to V$, $ae_0\mapsto av$, is surjective
and has kernel $\mathbb Qk$, where $k=e_0-e_1-z$.
Every nonidentity outside basis element annihilates both $f$
and $k$, while the identity fixes both. Hence $Af=\mathbb Qf$,
and $f\mapsto k$ is an injective $A$-linear map giving
\[
 0\longrightarrow Af\longrightarrow Ae_0
   \xrightarrow{\phi}V\longrightarrow0.
\]
Its first two terms are projective, proving the assertion.

The newly authored scout is
\path{monoid/code/scout_forced_three_row_actor_hulls13.py}.
The independent checker
\path{monoid/code/check_forced_three_row_actor_hulls13.py}
imports no scout or solver code. It reconstructs the complete
closed-set scope by closure-adjoining, rebuilds each literal
table and its full row action from band multiplication, and
checks every ideal, idempotent, basis, equivariance and exactness
equation above. Its receipt and the four full tables are under
\path{two_object19/ordinary13_forced_hulls/} in the registered
working directory. The scout's absent projective sections in
the last two cases are not themselves obstructions; the explicit
length-one presentations supply the proof.
\end{proof}

\begin{theorem}[Proved: three rows and seven outside elements]
\label{CP2:A-three-row-seven-outside-acts}
If a finite monoid has minimal band $3\times2$ and at most seven
outside elements, every finite ordinary right-act cone has zero
rational degree-two products. No bound on the state count or
homology is imposed.
\end{theorem}
\begin{proof}
If the monoid vertex has zero rational $H^2$, the relative
suspension argument already kills all second products. Otherwise
a nonzero balanced covector on the row augmentation module has
row values, up to affine change, either three distinct values
or $(0,0,1)$. With three distinct values, the nonconstant-column
actors fix a common rational affine row origin: identity-slope
actors fix every row, and a negative-slope actor is the unique
reflection. Lemma~\ref{CP2:A-act-common-affine-origin} applies.

In the two-value case a positive-slope row actor fixes row2 and
preserves $\{0,1\}$, whereas a negative-slope actor is $r_0$ or
$r_1$. If neither occurs, row2 is a common origin. If only one
occurs, every positive-slope actor fixes its selected row, since
composition would otherwise give the other negative-slope actor.
The midpoint of that selected row and row2 is then a common
origin. Thus a possible nonzero product forces both $r_i$ and
their squares $p_i$. This necessity argument is independent of
the outside-element budget.

Map the original monoid to its faithful action on the minimal
band. This is an actual surjective homomorphism, bijective on
that band. Its image has order at most thirteen and is one of
the four targets just classified, after relabelling. For an
original act $S$, use the actual induced act $S\otimes_M H$.
The integral comparison of
Theorem~\ref{CP2:A-band-quotient-integral-low-degree} gives an
isomorphism on rational second cohomology of the two act cones.
The target's row module has projective dimension at most one,
so Lemma~\ref{CP2:A-act-row-pd-one} kills all its second products.
Pullback gives the result. The target obstruction is unconditional;
neither boundedness nor projective-plane type is transferred through
the faithful quotient.
\end{proof}
